\chapter{HTTP API canónica}
\label{sec:api-http}

\section{Transporte y autenticación}

El Tauri daemon sirve la API en \code{http://127.0.0.1:\$GRAPHHUNTER_API_PORT}
(default \code{37891}, configurable por env var
\code{GRAPHHUNTER_API_PORT}). Por defecto sólo se bindea a
loopback; el host se puede cambiar con
\code{GRAPHHUNTER_API_HOST}.

Autenticación: cada request debe presentar el header
\code{Authorization: Bearer \$GRAPHHUNTER_API_TOKEN}. El token se
genera al instalar la app y se rotacional con el comando
\code{cli auth rotate}. Requests sin token reciben 401.

Formato: request y response en \code{application/json} UTF-8.
Timestamps ISO-8601 con offset explícito. Todos los identificadores
internos de entidad son enteros sin signo de 64 bits.

\section{Esquema introspectivo}

\begin{itemize}
    \item \code{GET /v1/schema} — Devuelve el listado completo de
    endpoints, sus DTOs y su estabilidad. Base para clientes
    generados automáticamente. Incluye además
    \code{schema_compute_time_ms} y \code{cache_hit} para
    diagnóstico.
\end{itemize}

\section{Salud y telemetría}

\begin{itemize}
    \item \code{GET /} — Root; info de build y versión.
    \item \code{GET /health} — Liveness / readiness combinado.
    \item \code{GET /metrics} — Contadores internos (hunts
    ejecutados, eventos ingeridos, errores por categoría). Formato
    compatible con Prometheus scrape.
\end{itemize}

\section{Catálogo y schema}

\begin{itemize}
    \item \code{GET /catalog} — Catálogo de tipos de entidad y
    relaciones del workspace.
    \item \code{GET /catalog/status} — Catálogo con flags de salud
    por entrada.
    \item \code{GET /catalog/diagnostics} — Resultados de
    verificación estructural.
    \item \code{GET /entity_types} — Lista de tipos de entidad con
    conteos.
    \item \code{GET /entity_type_counts} — Histograma de conteos.
    \item \code{GET /relation_schema} — Matriz
    \code{(from, relation, to)} observada en el grafo.
\end{itemize}

\section{Navegación del grafo}

\begin{itemize}
    \item \code{GET /search} — Búsqueda de entidades por valor o
    tipo. Parámetros: \code{q}, \code{type}, \code{limit}.
    \item \code{GET /expand} — Vecinos de un nodo.
    \item \code{GET /expand_grouped} — Vecinos agrupados por
    relation type (para UIs densas).
    \item \code{GET /node_details} — Propiedades y eventos asociados
    a un nodo.
    \item \code{POST /subgraph} — Subgrafo inducido por un conjunto
    de nodos.
    \item \code{GET /events_for_node} — Eventos donde participa un
    nodo.
    \item \code{GET /events_paginated} — Eventos con paginación
    (cursor-based, evita OOM sobre grafos grandes).
    \item \code{GET /path_nodes} — Nodos en un path concreto devuelto
    por un hunt.
    \item \code{GET /graph_summary} — Conteos agregados del grafo.
\end{itemize}

\section{Caza (hunts)}

\begin{itemize}
    \item \code{POST /run_hunt} — Ejecuta una hipótesis contra el
    grafo actual. Body: \code{\{dsl: string, min_score?: f32,
    max_results?: u32\}}. Devuelve
    \code{HuntResult\{paths, scores, explanation\}}.
    \item \code{GET /hunt_results} — Últimos resultados persistidos
    (paginados).
    \item \code{GET /explain_score} — Desglose de los cinco
    componentes de anomaly scoring para un path.
    \item \code{POST /diff_hunts} — Diferencia entre dos hunts por
    ID.
    \item \code{POST /parse_dsl} — Valida y devuelve AST de una
    hipótesis DSL sin ejecutar.
    \item \code{POST /compute_scores} — Recomputa scoring sobre
    resultados existentes (tras cambiar componentes).
    \item \code{POST /enable_anomaly_scoring} — Toggle de scoring
    anómalo a nivel sesión.
\end{itemize}

\section{Ingesta}

\begin{itemize}
    \item \code{POST /ingest_negotiate} — Acuerda con el cliente
    formato y mapping antes de abrir el stream.
    \item \code{POST /canonical_map} — Aplica un \code{FieldConfig}
    para normalizar un lote de eventos.
    \item \code{POST /parser_generate} — Deriva un \code{FieldConfig}
    heurístico a partir de un preview.
    \item \code{POST /schema_drift_detect} — Compara un batch contra
    el esquema conocido y reporta drift.
    \item \code{POST /mapping_regression_test} — Corre un mapping
    sobre un set congelado para detectar regresiones.
    \item \code{POST /publish_mapping} — Publica un mapping al
    \code{mapping_library}.
    \item \code{POST /fetch_shared_mappings} — Obtiene mappings
    compartidos del library.
\end{itemize}

\section{Anotación}

\begin{itemize}
    \item \code{POST /notes}, \code{POST /notes/} — Crea una nota
    sobre una entidad, path o hunt.
    \item \code{GET /notes/list} — Lista paginada de notas.
    \item \code{GET /notes/:id} — Nota puntual.
    \item \code{POST /tag} — Etiqueta una entidad.
    \item \code{POST /untag} — Remueve una etiqueta.
    \item \code{GET /tags} — Lista de entidades por etiqueta.
\end{itemize}

\section{Anomaly y enrichment}

\begin{itemize}
    \item \code{GET /temporal_heatmap} — Matriz hora $\times$ día de
    actividad.
    \item \code{GET /subnet_analysis} — Análisis por subnet
    (densidad, entidades top, eventos anómalos).
    \item \code{POST /enrich_ip} — Enriquecimiento online de una IP
    (GeoIP, WHOIS cached).
\end{itemize}

\section{Exportación}

\begin{itemize}
    \item \code{POST /export_iocs} — Exporta IOCs (STIX-like).
    \item \code{POST /publish_iocs_to_sentinel} — Publica IOCs a
    Azure Sentinel. Requiere \code{GRAPHHUNTER_IOC_PUBLISH_ENABLED=1}.
    \item \code{POST /save_hypothesis_as_sigma} — Convierte una
    hipótesis DSL a regla Sigma.
    \item \code{POST /file_ticket} — Abre ticket en el sistema
    configurado. Requiere \code{GRAPHHUNTER_TICKET_PUBLISH_ENABLED=1}.
    \item \code{POST /export_subgraph} — Exporta subgrafo en formato
    Graphson / JSONLines.
    \item \code{POST /export_ocsf} — Exporta eventos en OCSF v1.4.
    \item \code{POST /export_hunt_results} — Exporta resultados del
    último hunt en CSV/JSON.
\end{itemize}

\section{Invariantes}

\begin{itemize}
    \item \code{POST /check_invariants} — Corre checks estructurales
    sobre el grafo (treewidth, shape catalog, integridad temporal).
    \item \code{POST /invariant_check_hypothetical} — Corre los
    mismos checks como \emph{dry-run} sobre un grafo candidato
    (ingestión previa a commit).
\end{itemize}

\section{Agentic (slow lane)}

\begin{itemize}
    \item \code{POST /agentic_review_list} — Lista propuestas
    pendientes de review.
    \item \code{POST /agentic_review_approve} — Aprueba una
    propuesta.
    \item \code{POST /agentic_review_reject} — Rechaza una propuesta
    con razón.
\end{itemize}

\section{Dead Letter Queue}

\begin{itemize}
    \item \code{POST /list_dead_letters} — Lista eventos que
    fallaron ingesta.
    \item \code{POST /reingest_dead_letter} — Reintenta un evento
    (tras corregir mapping o parser).
    \item \code{POST /purge_dead_letters} — Borra entradas de la
    DLQ.
\end{itemize}

\section{SIEM (Sentinel)}

\begin{itemize}
    \item \code{POST /kql} — Ejecuta una query KQL contra Sentinel
    configurado y persiste resultados al workspace.
    \item \code{GET /sentinel_status} — Estado del conector y de la
    última sesión.
\end{itemize}

\section{Datasets y utilidad}

\begin{itemize}
    \item \code{GET /datasets} — Datasets disponibles en
    \filepath{demo_data/}.
\end{itemize}

\section{Códigos de respuesta}

\begin{itemize}
    \item \code{200} — éxito (con body JSON).
    \item \code{204} — éxito sin body (operaciones idempotentes).
    \item \code{400} — request inválido (DTO mal formado).
    \item \code{401} — token ausente o inválido.
    \item \code{403} — feature flag desactivado (p.ej.\ IOC publish
    sin env var).
    \item \code{404} — recurso inexistente (nodo, hunt, nota).
    \item \code{409} — conflicto (ingest en curso, session locked).
    \item \code{422} — validación semántica falló (DSL inválido,
    schema drift fuera de tolerancia).
    \item \code{500} — error interno; revisar logs del daemon.
\end{itemize}

\section{Total}

El daemon expone 64 rutas distintas (algunas con y sin slash
final). De ellas, 62 son endpoints de producto; \code{/} y
\code{/metrics} son meta. El listado \code{canónico} del runtime
(\code{GET /v1/schema}) es la fuente de verdad si hay discrepancia
con este documento.
